
\newif\ifcasscclass
\makeatletter
\@ifclassloaded{cas-sc}{\casscclasstrue}{\casscclassfalse}
\makeatother

\newif\ifpageisodd % È usato (nelle macro delle infopagini) dalla meccanica

\ifcasscclass
%\begingroup For image generation
  %\begingroup Tipografia
    \newcommand\textsize[1]{%
      \pgfmathparse{#1*1.2}%
      \fontsize{#1pt}{\pgfmathresult pt}\selectfont%
    }

    \DeclareMathOperator\FRC{\mathcal R}

    % (il rilevamento della classe è in cima al file)

    \providecommand\Ifthispageodd[2]{\checkoddpage\ifoddpage#1\else#2\fi}
  %\endgroup
\fi

  %\begingroup TikZ
    %\begingroup Didascalie
      \newcommand\newcaption[2]{%
        \captionsetup{width=#2} % aboveskip=-.2cm,
        \caption{#1}%
      }
      
      \newcommand\newsubcaption[2][ ]{%
        \refstepcounter{subfigure}%
        (\thesubfigure)#1#2%
      }
    %\endgroup
    
\ifcasscclass
    %\begingroup RedefineTikZbounds
      % La differenza è che colla prima bisogna catturare anche le etichette perché si crea un pdf isolato mentre colla non è importante perché viene tutto compilato sul momento

      %\begingroup Versione valida con esternalizzazione [ma anche senza]
        \newcommand\redefineTikZbounds[2]{
          \coordinate (RealSW) at (current bounding box.south west);
          \coordinate (RealNE) at (current bounding box.north east);
          \coordinate (Center) at ($(#1.west)!0.5!(#2.east)$);
    
          \pgfresetboundingbox 

          \useasboundingbox 
            let \p1=(RealSW),\p2=(RealNE),\p3=(Center),
            \n1={max(\x3-\x1,\x2-\x3)} in
            (\x3-\n1,\y1) rectangle (\x3+\n1,\y2+.25em);

          % \path[use as bounding box] 
          %   let \p1=(RealSW),\p2=(RealNE),\p3=(Center),
          %   \n1={max(\x3-\x1,\x2-\x3)} in
          %   (\x3-\n1,\y1) rectangle (\x3+\n1,\y2+.25em);
          % Lo «.25em» discende dalla traslazione imposta all'etichetta della scala verticale mediante «yScaleLabelShift»
        } % Scatola delimitante simmetrica
      %\endgroup

      %\begingroup Versione valida senza esternalizzazione
        % \newcommand\redefineTikZbounds[2]{
        %   % Salva le coordinate della cima e del fondo della figura corrente (cosí da mantenere i titoli e le didascalie)
        %   \coordinate (Cima) at (current bounding box.north);
        %   \coordinate (Fondo) at (current bounding box.south);
    
        %   \pgfresetboundingbox % Resetta la scatola dei limiti
    
        %   \useasboundingbox % Definisce la nuova scatola come un rettangolo
        %     (#1.west |- Cima)  % con ascissa sinistra data da «Col1 c1r1.west», 
        %     rectangle      % ascissa destra data da «Col2 c1r1.east»
        %     (#2.east |- Fondo); % ordinate date dalla cima e fondo originali
        %     % L'operatore «|-» in «(A |- B)» prende come ascissa quella del punto A mentre come coordinata quella del punto B
        % }
      %\endgroup
    %\endgroup

    %\begingroup TikZBox
      % Parzialmente ispirato da https://gist.github.com/Caramdir/899107

      \makeatletter
      \newsavebox\tikzBufferBox
      \xdef\offsetX{0pt}
      \xdef\offsetY{0pt}
      
      % \tikzbox[#1]{#2}{#3}[#4]{#5}
      % #1 -> Ancora rispetto cui posizionare l'intera scatola (di base è «south west», ossia l'angolo in basso a sinistra)
      % #2 -> Nodo dove posizionare l'origine (per la natura di come LaTeX riceve gli argomenti è buon pratica passare questo circondato da parentesi graffe)
      % #3 -> Nome della scatola limitante definita dall'ambiente «scope»
      % #4 -> Argomenti opzionali da passare all'ambiente «scope»; di base sono vuoti
      % #5 -> Contenuto della scatola circondato dall'ambiente «scope»
      % Es. \tikzbox{{(0,0)}}{figuraTikZ}{blocco TikZ}

      % Macro di apertura che controlla se v'è l'argomento facoltativo dell'ancora
      \def\tikzbox{\pgfutil@ifnextchar[{\tikzbox@anchor}{\tikzbox@anchor[south west]}}
      % La funzione «\pgfutil@ifnextchar» serve per controllare se il prossimo carattere è l'apertura di una paretensi graffa «[»
      
      % Macro intermedia per cogliere gli argomenti obbligatori e controllare quelli successivi facoltativi
      \def\tikzbox@anchor[#1]#2#3{\pgfutil@ifnextchar[{\tikzbox@opt[#1]#2#3}{\tikzbox@opt[#1]#2#3[]}}

      % Macro degli argomenti obbligatori
      \long\def\tikzbox@opt[#1]#2#3[#4]#5{%
        \node[ % Misurazione della scatola in un nodo invisibile
          overlay,%
          opacity=0,%
          inner sep=0pt,%
          outer sep=0pt%
        ] {%
          \sbox{\tikzBufferBox}{%
            \begin{tikzpicture}%
              \begin{scope}[local bounding box=measuringTikZBox]%
                #5%
              \end{scope}%
              \pgfpointanchor{measuringTikZBox}{#1}%
              \pgfgetlastxy{\tempX}{\tempY}% 
              \xdef\offsetX{\tempX}%
              \xdef\offsetY{\tempY}%
            \end{tikzpicture}%
          }%
        };%
        %
        % Disegno reale traslato
        \begin{scope}[%
          local bounding box=#3,%
          shift={#2},%
          xshift=-\offsetX,%
          yshift=-\offsetY,%
          #4%
        ]%
          #5%
        \end{scope}%
      }%
      \makeatother
    %\endgroup

    %\begingroup Chiavi-Valori
      %\begingroup Testo
        % La dimensione del testo di base è 11pt
        \newcommand\titleSize{11}
        \newcommand\subCaptionSize{10}
        \newcommand\ParameterSize{9}
        \newcommand\labelSize{10}
        \newcommand\tickSize{9}
        \newcommand\legendSize{9}
      %\endgroup
    
      %\begingroup Infopagine
        %\begingroup Chiavi condivise
        \newcounter{lpc} % «loopcounter»
        \newcommand\infopageSharedKeys{
          %\begingroup Testo
            \renewcommand\titleSize{10}
            \renewcommand\subCaptionSize8
            \renewcommand\ParameterSize7
            \renewcommand\labelSize9
            \renewcommand\tickSize6
            \renewcommand\legendSize7
          %\endgroup
    
          %\begingroup Dimensioni
            \pgfmathsetmacro\plotWidth7
            \pgfmathsetmacro\plotHeight{.618*\plotWidth}
            \pgfmathsetmacro\halfPlotHeight{\plotHeight/2}
            %
            \pgfmathparse{0.825*\plotWidth}
            \pgfmathsetmacro\plotWidth\pgfmathresult
            \pgfmathsetmacro\halfWidth{\plotWidth/2}
            \pgfmathsetmacro\plotLastWidth{1.25*\plotWidth}
    
            \pgfmathsetmacro\xPlotSep{1}
            \pgfmathsetmacro\yPlotSep{1.25}
            \pgfmathsetmacro\yGroupSep{2.75}
          %\endgroup
    
          %\begingroup Legende
            \newcommand\placeLegend[5]{%
              /pgfplots/plotLegend,
              at={([shift={(##4,##5)}]axis description cs:##1,##2)},
              anchor=##3,
            }
            %
            \pgfplotsset{%
              plotLegend/.style={%
                font=\textsize{\legendSize},
                /pgfplots/legend cell align=left,
                % align=left,
                % draw=white!15!black,
                draw=none, % Bordo invisibile
                fill=white,
                fill opacity=0.5,
                text opacity=1,
                % text width=2cm,
              },
              plotLegendNW/.style={legend style={\placeLegend{0}{1}{north west}{0em}{0em}}},
              plotLegendNE/.style={legend style={\placeLegend{1}{1}{north east}{0em}{0em}}},
              plotLegendSW/.style={legend style={\placeLegend{0}{0}{south west}{0em}{0em}}},
              plotLegendSE/.style={legend style={\placeLegend{1}{0}{south east}{0em}{0em}}},
              %
              % plotLegendW/.style={legend style={plotLegend,at={(0,.5)},anchor=west}},
              %
              % plotLegendNW/.style={legend style={plotLegend,legend pos=north west}},
              % plotLegendNE/.style={legend style={plotLegend,legend pos=north east}},
              % plotLegendSW/.style={legend style={plotLegend,legend pos=south west}},
              % plotLegendSE/.style={legend style={plotLegend,legend pos=south east}},
              %
              % plotLegendNWcloser/.style={legend style={plotLegend,at={(0,1)},anchor=north west}},
              % plotLegendNEcloser/.style={legend style={plotLegend,at={(1,1)},anchor=north east}},
              % plotLegendSWcloser/.style={legend style={plotLegend,at={(0,0)},anchor=south west}},
              % plotLegendSEcloser/.style={legend style={plotLegend,at={(0,1)},anchor=south east}},
            }
            %
            \pgfmathsetmacro\lineLegendLength{1}
            \pgfmathsetmacro\areaLegendHeigth{.25}
            \pgfplotsset{
              legend image code/.code={
                \draw [mark repeat=2,mark phase=2,##1]
                  plot coordinates {
                    (0em,0em)
                    (.5*\lineLegendLength em,0em)
                    (\lineLegendLength em,0em)
                  };
              },
              /pgfplots/area legend/.style={
                legend image code/.code={
                  \draw [##1,fill,draw=none] 
                    (0em,-\areaLegendHeigth em) 
                    rectangle 
                    (\lineLegendLength em,\areaLegendHeigth em);
                },
              },
            }
            %
            \newcommand\ExactHistogramLegend{Exact histogram}
            \newcommand\ApproxHistogramLegend{Approx. histogram}
            \newcommand\RealHistogramLegend{Real histogram}
            %
            \newcommand\ExactBLNFitLegend{Exact BLN fit}
            \newcommand\ApproxBLNFitLegend{Approx BLN fit}
            \newcommand\RealBLNFitLegend{Real BLN fit}
            %
            \newcommand\ExactEmpiricCCDFLegend{Exact empiric CCDF}
            \newcommand\ApproxEmpiricCCDFLegend{Approx. empiric CCDF}
            \newcommand\RealEmpiricCCDFLegend{Real empiric CCDF}
            %
            \newcommand\ExactParetoFitLegend{Exact Pareto fit}
            \newcommand\ApproxParetoFitLegend{Approx. Pareto fit}
            \newcommand\RealParetoFitLegend{Real Pareto fit}
            %
            \newcommand\ExactScatterLegend{Exact scatter}
            \newcommand\ApproxScatterLegend{Approx. scatter}
            \newcommand\RealScatterLegend{Real scatter}
            %
            \newcommand\ClasskILegend{Class k=283}
            \newcommand\ClasskIILegend{Class k=83}
            \newcommand\ClasskIIILegend{Class k=43}
            \newcommand\ClasskIVLegend{Class k=10}
          %\endgroup
    
          %\begingroup Etichette vuote
            \pgfplotsset{
              yLabelEmpty/.style={ylabel={}},
              yTickLabelsEmpty/.style={yticklabels={}},
              yTickScaleLabelEmpty/.style={ytick scale label code/.code={}},
              xTickScaleLabelEmpty/.style={xtick scale label code/.code={}},
              emptyLeged/.style={
                legend style={draw=none,fill=none,nodes={draw=none}},
                legend image post style={opacity=0},
              },
            }
          %\endgroup
    
          %\begingroup Estetica
            \newcommand\smallminus{\scalebox{0.5}[1]{$-$}\hspace{-.11em}}
            \pgfplotsset{
              /pgfplots/log number format basis/.code 2 args={%
                \pgfmathsetmacro\exponent{##2}%
                \pgfmathparse{sign(\exponent)}%
                \ifnum\pgfmathresult<0%
                  \pgfmathparse{abs(\exponent)}%
                  $##1^{\smash{\smallminus\pgfmathprintnumber{\pgfmathresult}}}$%
                \else%
                  $##1^{\smash{\pgfmathprintnumber{\exponent}}}$%
                \fi%
              },
              ytick scale label code/.code={%
                \pgfmathsetmacro\exponent{##1}%
                \pgfmathparse{sign(\exponent)}%
                \ifnum\pgfmathresult<0
                  \pgfmathparse{abs(\exponent)}%
                  $\cdot 10^{\smash{\smallminus\pgfmathprintnumber{\pgfmathresult}}}$%
                \else
                  $\cdot 10^{\smash{\pgfmathprintnumber{\exponent}}}$%
                \fi
              },
            }
          %\endgroup
    
          %\begingroup Colori
            \definecolor{blueish}{rgb}{0.00000,0.44700,0.74100}%
            \definecolor{reddish}{rgb}{0.83529,0.36863,0.00000}%
          %\endgroup
    
          %\begingroup Commenti
            %\begingroup Impostazioni
            \tikzset{
              title/.style={
                align=center,
                anchor=south,
                font=\bfseries\textsize{\titleSize},
              },
              fittedParameters/.style={
                align=center,
                font=\textsize{\ParameterSize},
              },
              varyingParameter/.style={
                align=center,
                % rotate=90,
                font=\textsize{\subCaptionSize},
              },
              subCaption/.style={
                % text width=\plotWidth cm,
                align=center,
                font=\textsize{\subCaptionSize},
              },
            }
            \newcommand\codetocolor[1]{\ifcase##1blueish\or reddish\else black\fi}
            %\endgroup
    
            %\begingroup Parametri Sardegna
            \newcommand\betaReal{\num{1.273}}
            \newcommand\zetaReal{\num{0.929}}
            \newcommand\mIReal{\num{3.227}}
            \newcommand\sIReal{\num{0.408}}
            \newcommand\mIIReal{\num{4.099}}
            \newcommand\sIIReal{\num{0.556}}
            %\endgroup
    
            %\begingroup Parametri Pareto
            \newcommand\codetoParetovalue[1]{
              \ifcase##1\betaExact\or\betaAprx\else\betaReal\fi
            }
            \newcommand\writeParetoParameters[7][center]{
              \node[fittedParameters,anchor=##1] at ([shift={(##5,##6)}]##4) {%
                \(\begin{aligned}
                  \color{\codetocolor{##2}}\beta
                  &\color{\codetocolor{##2}}=\codetoParetovalue{##2}\\[##7]
                  \color{\codetocolor{##3}}\beta
                  &\color{\codetocolor{##3}}=\codetoParetovalue{##3}
                \end{aligned}\)
              };
            }
            %\endgroup

            %\begingroup % Vecchia versione in cui si scrivevano esplicitamente tutte le 9 combinazioni possibili
            % \newcommand\writeParameters[6]{%
            %   \node[parameter] at ([shift={(##4,##5)}]##1) {%
            %     \ifcase##2% Caso 0
            %       \ifcase##3% Caso 0
            %         \(\begin{aligned}
            %           \color{blueish}\beta&\color{blueish}=\betaExact\\[##6]
            %           \color{blueish}\beta&\color{blueish}=\betaExact
            %         \end{aligned}\)
            %       \or% Caso 1
            %         \(\begin{aligned}
            %           \color{blueish}\beta&\color{blueish}=\betaExact\\[##6]
            %           \color{reddish}\beta&\color{reddish}=\betaAprx% Caso 1
            %         \end{aligned}\)
            %       \else% Caso 2
            %         \(\begin{aligned}
            %           \color{blueish}\beta&\color{blueish}=\betaExact\\[##6]
            %           \color{black}\beta&\color{black}=\betaReal% Caso 2
            %         \end{aligned}\)
            %       \fi
            %     \or% Caso 1
            %       \ifcase##3% Caso 0
            %         \(\begin{aligned}
            %           \color{reddish}\beta&\color{reddish}=\betaAprx\\[##6]
            %           \color{blueish}\beta&\color{blueish}=\betaExact
            %         \end{aligned}\)
            %       \or% Caso 1
            %         \(\begin{aligned}
            %           \color{reddish}\beta&\color{reddish}=\betaAprx\\[##6]
            %           \color{reddish}\beta&\color{reddish}=\betaAprx
            %         \end{aligned}\)
            %       \else% Caso 2
            %         \(\begin{aligned}
            %           \color{reddish}\beta&\color{reddish}=\betaAprx\\[##6]
            %           \color{black}\beta&\color{black}=\betaReal
            %         \end{aligned}\)
            %       \fi
            %     \else% Caso 2
            %       \ifcase##3% Caso 0
            %         \(\begin{aligned}
            %           \color{black}\beta&\color{black}=\betaReal\\[##6]
            %           \color{blueish}\beta&\color{blueish}=\betaExact
            %         \end{aligned}\)
            %       \or% Caso 1
            %         \(\begin{aligned}
            %           \color{black}\beta&\color{black}=\betaReal\\[##6]
            %           \color{reddish}\beta&\color{reddish}=\betaAprx
            %         \end{aligned}\)
            %       \else% Caso 2
            %         \(\begin{aligned}
            %           \color{black}\beta&\color{black}=\betaReal\\[##6]
            %           \color{black}\beta&\color{black}=\betaReal
            %         \end{aligned}\)
            %       \fi
            %     \fi
            %   };
            % }
            %\endgroup
    
            %\begingroup Parametri bilognormale
            \newcommand\codetoparameter[1]{
              \ifcase##1\zeta  % Caso 0
              \or\mu_1     % Caso 1
              \or\sigma_1   % Caso 2
              \or\mu_2     % Caso 3
              \else\sigma_2\fi % Caso 4
            }
            \newcommand\codetoBiLognormalvalue[2]{
              \ifcase##1 % Caso 0
                \ifcase##2\zetaExact % Caso 0
                \or\mIExact    % Caso 1
                \or\sIExact    % Caso 2
                \or\mIIExact    % Caso 3
                \else\sIIExact\fi % Caso 4
              \or % Caso 1
                \ifcase##2\zetaAprx % Caso 0
                \or\mIAprx     % Caso 1
                \or\sIAprx     % Caso 2
                \or\mIIAprx    % Caso 3
                \else\sIIAprx\fi  % Caso 4
              \else % Caso 2
                \ifcase##2\zetaReal % Caso 0
                \or\mIReal     % Caso 1
                \or\sIReal     % Caso 2
                \or\mIIReal    % Caso 3
                \else\sIIReal\fi  % Caso 4
              \fi
            }
            \newcommand\matrixline[2]{\codetoparameter{##2}&=\codetoBiLognormalvalue{##1}{##2}}
            \def\lineloop##1##2{%
              \matrixline{##1}{\number\value{lpc}}\stepcounter{lpc}%
              \ifnum\number\value{lpc}<5\\[##2]\lineloop{##1}{##2}\fi%
            } % https://tex.stackexchange.com/questions/491358/tikz-foreach-to-work-with-the-aligned-environment
            \newcommand\writeBiLognormalParameters[6][center]{
              \node[fittedParameters,anchor=##1] at ([shift={(##4,##5)}]##3) {%
                \setcounter{lpc}{0}
                \(\color{\codetocolor{##2}}
                \begin{aligned}
                  \lineloop{##2}{##6}
                \end{aligned}\)
              };
            }
            %\endgroup
    
            %\begingroup Altre etichette
            \newcommand\writeSizeEvolutionDescriptions[1][5]{%
              \foreach \col/\clr/\txt in {%
                1/blueish/Esatto,
                2/reddish/Approssimato
              }{
                \node[
                  fittedParameters,
                  anchor=south,
                  yshift=.1em,
                ] at (Row##1 c\col r1.south) {\color{\clr}\txt};
              }
            }
    
            \newcommand\writeVaryingParameter[5]{%
              \foreach \r in {1,...,5}{
                \node[
                  varyingParameter,
                  anchor=##3,
                ] at ([shift={(##4,##5)}]Row\r\space c##1r1.##2)
                {\(\varyingParameterSymbol=\varyingParameterList\r\)};
              }
            }
            %\endgroup
          %\endgroup
    
          %\begingroup Sottodidascalie
            \newcommand\writeSubCaption[5]{%
              \node[subCaption,anchor=##3]
              at ([shift={(##4,##5)}]##1.##2)
              {\newsubcaption[]{}};
            }
            \newcommand\writeSubCaptions[5][3]{%
              \foreach \r in {1,...,5}{
                \foreach \c in {1,...,##1}{
                  \writeSubCaption{Row\r\space c\c r1}{##2}{##3}{##4}{##5}
                }
              }
            }
          %\endgroup
        }
        %\endgroup
    
        %\begingroup Chiavi del primo tipo d'infopagina
        \newcommand\infopageKeysI{
          \infopageSharedKeys
    
          %\begingroup Gruppi
            \ifpageisodd
              \pgfplotsset{/pgfplots/group/yDescPos/.style={y descriptions at=edge right}}
            \else
              \pgfplotsset{/pgfplots/group/yDescPos/.style={y descriptions at=edge left}}
            \fi
            \pgfplotsset{
              /pgfplots/group/plotGroup/.style={
                horizontal sep=\xPlotSep em,
                vertical sep=\yPlotSep em,
                yDescPos,
                group size=3 by 1,
              },
              /pgfplots/group/plotLastGroup/.style={
                plotGroup,
                x descriptions at=edge bottom,
              },
            }
          %\endgroup
    
          %\begingroup Stili
            %\begingroup Selezione «recto»/«verso»
            \ifpageisodd
              \pgfplotsset{%
                yLineAxis/.style={axis y line*=right},
                /tikz/yLabelShift/.style={%
                  yshift=-1.25em,
                  at={(axis description cs:1,0.5)},
                  anchor=north,
                  % rotate=180,
                },
                /tikz/yScaleLabelShift/.style={%
                  yshift=.25em,
                  at={(1,1)},
                  anchor=south east,
                },
                /tikz/yTickLabelShift/.style={xshift=-0.15em},
              }
            \else
              \pgfplotsset{%
                yLineAxis/.style={axis y line*=left},
                /tikz/yLabelShift/.style={%
                  yshift=1.25em,
                  at={(axis description cs:0,0.5)},
                  anchor=south,
                  % rotate=180,
                },
                /tikz/yScaleLabelShift/.style={%
                  yshift=.25em,
                  at={(0,1)},
                  anchor=south west,
                },
                /tikz/yTickLabelShift/.style={xshift=+0.15em},
              }
            \fi
            %\endgroup
            %
            %\begingroup Grafici
            \pgfplotsset{
              plotRow/.style={
                scale only axis,
                width=\plotWidth cm,height=\plotHeight cm,
                xlabel style={font=\color{white!15!black}\textsize{\labelSize}},
                ylabel style={font=\color{white!15!black}\textsize{\labelSize},yLabelShift},
                xticklabel style={font=\textsize{\tickSize}},
                yticklabel style={font=\textsize{\tickSize},yTickLabelShift},
                y tick scale label style={yScaleLabelShift},
                % axis x line*=bottom,yLineAxis,
                enlarge x limits=0.1,enlarge y limits=0.1,
                xmajorgrids,ymajorgrids,
                xminorticks=true,yminorticks=true,
                scale ticks above exponent=2,
                axis background/.style={fill=white},
                grid style={dashed}
              },
              % xlabel123/.style={xlabel={\(h\)}},
              xlabel3/.style={xlabel={\(h\)}},
              ylabel23/.style={ylabel={\(\FRC(h)\)}},
              row1Keys/.style={
                % ylabel={\(\underhat{\bar l}_N(h)\)},
                ylabel={\(\bar l_N(h)\)},
                % enlarge x limits={lower,value=.05},
                enlarge x limits=.05,
                enlarge y limits={upper,value=.05},
                xmode=linear,
                log origin=infty,
              },
              row2KeysWithLabels/.style={
                ylabel23,
                xmode=linear,ymode=log
              },
              row2Keys/.style={
                ylabel23,
                xmode=linear,ymode=log,
                xticklabels={},
              },
              row3Keys/.style={
                xlabel3,ylabel23,
                xmode=linear,ymode=log
              },
              row4Keys/.style={
                xlabel={\(\xi\)},ylabel={\(s(\xi)\)},
                xmode=log,ymode=log
              },
              row5KeysLin/.style={
                width=\plotLastWidth cm,
                xlabel={\(t\)},ylabel={\(s(t)\)},
                enlarge x limits=0,
                enlarge y limits=0.05,
                x tick scale label style={
                  yshift=-1em,
                  at={(.95,0)},
                  anchor=north,
                },
              },
              row5Keys/.style={
                row5KeysLin,
                ymode=log,
              },
            }
            %\endgroup
          %\endgroup
        }
        %\endgroup
    
        %\begingroup Chiavi del secondo tipo d'infopagina
        \newcommand\infopageKeysII{
          \infopageSharedKeys
    
          %\begingroup Gruppi
            \ifpageisodd
              \pgfplotsset{/pgfplots/group/yDescPos/.style={y descriptions at=edge right}}
            \else
              \pgfplotsset{/pgfplots/group/yDescPos/.style={y descriptions at=edge left}}
            \fi
            \pgfplotsset{
              /pgfplots/group/plotGroup/.style={
                horizontal sep=\xPlotSep em,
                vertical sep=\yPlotSep em,
                y descriptions at=all,
                group size=3 by 1,
              },
              /pgfplots/group/plotLastGroup/.style={
                plotGroup,
                x descriptions at=edge bottom,
              },
            }
          %\endgroup
    
          %\begingroup Stili
            %\begingroup Selezione «recto»/«verso»
            \ifpageisodd
              \pgfplotsset{
                yLabel/.style={
                  ylabel style={
                    font=\color{white!15!black}\textsize{\labelSize},
                    yshift=-1.25em,
                    at={(axis description cs:1,0.5)},
                    anchor=north,
                    % rotate=180,
                  },
                },
                yLabelInv/.style={yLabelEmpty},
                yScaleLabelShift/.style={
                  y tick scale label style={
                    yshift=.25em,
                    % xshift=1em,
                    at={(1,1)},
                    anchor=south east,
                  },
                },
                yScaleLabelShiftInv/.style={
                  y tick scale label style={
                    yshift=.25em,
                    % xshift=-1em,
                    at={(0,1)},
                    anchor=south west,
                  },
                },
                % yLineAxis/.style={axis y line*=right},
              }
            \else
              \pgfplotsset{
                yLabel/.style={yLabelEmpty},
                yLabelInv/.style={
                  ylabel style={
                    font=\color{white!15!black}\textsize{\labelSize},
                    yshift=1.25em,
                    at={(axis description cs:0,0.5)},
                    anchor=south,
                    % rotate=180,
                  },
                },
                yScaleLabelShift/.style={
                  y tick scale label style={
                    yshift=.25em,
                    % xshift=1em,
                    at={(1,1)},
                    anchor=south east,
                  },
                },
                yScaleLabelShiftInv/.style={
                  y tick scale label style={
                    yshift=.25em,
                    % xshift=-1em,
                    at={(0,1)},
                    anchor=south west,
                  },
                },
                % yLineAxis/.style={axis y line*=left},
              }
            \fi
            %\endgroup
            %
            %\begingroup Grafici
            \pgfplotsset{
              plotRow/.style={
                scale only axis,
                width=\plotWidth cm,height=\plotHeight cm,
                xlabel style={font=\color{white!15!black}\textsize{\labelSize}},
                % ylabel style={font=\color{white!15!black}\textsize{\labelSize},yLabelShift},
                xticklabel style={font=\textsize{\tickSize}},
                % axis x line*=bottom,yLineAxis,
                enlarge x limits=0.1,enlarge y limits=0.1,
                xmajorgrids,ymajorgrids,
                xminorticks=true,yminorticks=true,
                scale ticks above exponent=2,
                axis background/.style={fill=white},
                grid style={dashed},
              },
              plotLeftCol/.style={
                yLabelInv,
                yticklabel style={font=\textsize{\tickSize},xshift=0.15em},
                yticklabel pos=left,
                yScaleLabelShiftInv,
              },
              plotMiddleCol/.style={
                yLabelEmpty,
                yTickLabelsEmpty,
                yTickScaleLabelEmpty,
              },
              plotRightCol/.style={
                yLabel,
                yticklabel style={font=\textsize{\tickSize},xshift=-0.15em},
                yticklabel pos=right,
                yScaleLabelShift,
              },
              xlabel3/.style={xlabel={\(h\)}},
              ylabel23/.style={ylabel={\(\FRC(h)\)}},
              row1Keys/.style={
                ylabel={\(\bar l_N(h)\)},
                enlarge x limits=.05,
                enlarge y limits={upper,value=.05},
                xmode=linear,
                log origin=infty,
              },
              row2Keys/.style={
                ylabel23,
                xmode=linear,ymode=log,
                % xticklabels={},
              },
              row3Keys/.style={
                xlabel3,ylabel23,
                xmode=linear,ymode=log,
              },
              row4KeysLin/.style={
                xlabel={\(\xi\)},ylabel={\(s(\xi)\)},
                xmode=log,
              },
              row4Keys/.style={row4KeysLin,ymode=log},
              row5KeysLin/.style={
                width=\plotLastWidth cm,
                xlabel={\(t\)},ylabel={\(s(t)\)},
                enlarge x limits=0,
                enlarge y limits=0.05,
                x tick scale label style={
                  yshift=-1em,
                  at={(.95,0)},
                  anchor=north,
                },
              },
              row5Keys/.style={
                row5KeysLin,
                ymode=log,
              },
            }
            %\endgroup
          %\endgroup
        }
        %\endgroup
      %\endgroup
    %\endgroup

    %\begingroup Diagrammi dei nodi
    \newcommand\nodeDiagramSettings{
      %\begingroup Dimensioni
        \pgfmathsetmacro\xPlotSep{10}
        \pgfmathsetmacro\yPlotSep{4}
        \pgfmathsetmacro\plotScale{.8}

        \pgfmathsetmacro\titleSize{11}
        \pgfmathsetmacro\commentSize{10}

        \pgfmathsetmacro\nodeSize{20pt} % Diametro del nodo
        \pgfmathsetmacro\starSize{2.25*\nodeSize}
        \pgfmathsetmacro\numberSize{.53*\nodeSize}
      %\endgroup

      %\begingroup Colori
        % \definecolor{cRed}{HTML}{D68490}
        % \definecolor{cPurple}{HTML}{7E6EAC}
        % \definecolor{cTeal}{HTML}{87BDB1}
        % \definecolor{cOrange}{HTML}{F2D0A4}
        % \definecolor{cEdgeLight}{HTML}{EEE8D5} 
        % \definecolor{cEdgeDark}{HTML}{998888}
        % \definecolor{cStar}{HTML}{FDF6E3}
        %
        \definecolor{color1}{HTML}{800000} % Deep Red
        \definecolor{color2}{HTML}{4682B4} % Steel Blue
        \definecolor{color3}{HTML}{6B8E23} % Olive Green
        \definecolor{color4}{HTML}{708090} % Slate Gray
        \definecolor{color5}{HTML}{BC8F8F} % Rosy Brown
        \definecolor{color6}{HTML}{008B8B} % Dark Cyan
        \definecolor{color7}{HTML}{B8860B} % Dark Goldenrod
      %\endgroup

      %\begingroup Testo
        \tikzset{
          titleStyle/.style={
            font=\textsize{\titleSize},
            align=center,
            anchor=south,
          },
          commentStyle/.style={font=\textsize{\commentSize}},
          arrowCommentStyle/.style={
            midway,
            commentStyle,
            left,
            sloped,
            below
          },
          numberStyle/.style={font=\textsize{\numberSize}}
        }
      %\endgroup

      %\begingroup Frecce
        \tikzset{
          transformingArrow/.style={
            % =Latex,
            ->,
            >=Stealth,
            dashed,
            line width=.8pt,
            black!90,
            shorten >=5pt,
            shorten <=10pt
          },
          indirectArrow/.style={
            draw=black!95,
            line width=.5pt,
          },
          nodeArrowExact/.style={
            indirectArrow,
            % ->,
            % -{Stealth[scale=1.2]},
            -{Triangle[scale=.75]},
            % draw=cEdgeLight,
            % draw=red!50,
            % shorten >=2pt,
            % shorten <=2pt
          },
          nodeArrowApprox/.style={
            indirectArrow,
            -,
            dashed,
          },
          transparentNodeArrow/.style={
            opacity=.4,
          },
          arrowAxis/.style={
            ->,
            line width=.3pt,
            >=Latex,
          },
        }
      %\endgroup

      %\begingroup Nodi
        \tikzset{
          nodeStyle/.style={
            circle,
            minimum size=\nodeSize,
            % draw=black!80,
            draw=black,
            line width=.5pt,
            inner sep=0pt,
          },
          starStyle/.style={
            star,
            star points=5,
            star point ratio=2,
            % fill=cStar,
            fill=black!50,
            fill opacity=.25,
            minimum size=\starSize,
            draw=none
          },
        }
      %\endgroup
    }
    %\endgroup

    %\begingroup GNU-Plot
      % \IfFileExists{\sourceauxpath.aux}{
      %   \externaldocument{\sourceauxpath}
      % }{\relax}
      
      \pgfkeys{/pgf/plot/gnuplot call={cd /d "\sourceoutdir" && gnuplot}}
    %\endgroup
  %\endgroup
\fi
%\endgroup

%\begingroup For the paper
  \def\tsc#1{\csdef{#1}{\textsc{\lowercase{#1}}\xspace}}
  \tsc{WGM}
  \tsc{QE}
  
  \newcommand\etal{et al.}
  \newcommand\KTMAS{\hyperlink{acrKineticTheoryMultiAgentSystems}{KTMAS}}
  
  % \newcommand\MathSpacing[3]{
  %   \thinmuskip = #1mu\relax
  %   \medmuskip = #2mu\relax
  %   \thickmuskip = #3mu\relax
  % }
  % \MathSpacing123
  
  \newcommand\RPlus{\mathbb R_+}
  \newcommand\NPlus{\mathbb N_+}
  \DeclareMathAlphabet\pazocal{OMS}{zplm}{m}{n}
  
  \DeclareMathOperator\Lognormal{Lognormal}
  \DeclareMathOperator\BiLognormal{BiLognormal}
  \DeclareMathOperator\Bernoulli{Bernoulli}
  \DeclareMathOperator\Exponential{Exp}
  \DeclareMathOperator\Student{Student}
  \DeclareMathOperator\Variance{var}
  \DeclareMathOperator\IC{IC}
  \DeclareMathOperator\Histogram{\mathcal B}
  \DeclareMathOperator\CCDF{\mathcal R}
    
  \newcommand\ThreeQuarters{{3\hspace{-0.15em}/\hspace{-0.15em}4}}
  \newcommand\OneQuarter{{1\hspace{-0.15em}/\hspace{-0.15em}4}}
  
  \newcommand\de[1][]{\mathrm{d}#1}
  \newcommand\derS[3][d]{\frac{\mathrm{#1}#2}{\mathrm{#1}#3}}
  \newcommand\derP[2]{\frac{\partial#1}{\partial#2}}
  \newcommand\underhat[1]{\underaccent{\hat}{#1}}
  \renewcommand\undertilde[1]{\underaccent{\tilde}{#1}}
  
  \makeatletter
  \newcommand\eqtag{\@ifstar{\eqtag@star}{\eqtag@nostar}}
  
  \newcounter{eqtag}
  \newcommand\eqtag@nostar[2]{%
    \begin{equation*}%
      \refstepcounter{eqtag}%
      % \protected@edef\@currentlabeltype{eqtag}%
      \protected@edef\@currentlabel{#1}%
      #2%
      \tag*{#1}%
    \end{equation*}%
  }
  \newcommand\eqtag@star[2]{%
    \begin{equation*}%
      \refstepcounter{eqtag}%
      % \protected@edef\@currentlabeltype{eqtag}%
      \protected@edef\@currentlabel{#1}%
      \let\label\ltx@label%
      #2%
      % \notag%
    \end{equation*}%
  }
  \makeatother
  
  \newcommand\crefsubfigInt[4][-]{Figures~\ref{#2}#3#1\ref{#2}#4}
  \newcommand\crefsubfigSep[4]{Figures~\ref{#1}#2\foreach \l in#3{,~\ref{#1}\l}~and~\ref{#1}#4}
  
  \crefname{section}{§}{§§}
  \crefrangeformat{section}{§§~#3#1#4--#5#2#6}
  \crefmultiformat{section}{§§~#2#1#3}{ and~#2#1#3}{, #2#1#3}{ and~#2#1#3}
  \crefrangemultiformat{section}{§§~#3#1#4--#5#2#6}{ and~#3#1#4--#5#2#6}{, #3#1#4--#5#2#6}{ and~#3#1#4--#5#2#6}
  
  \crefname{equation}{}{}
  \crefrangeformat{equation}{(#3#1#4--#5#2#6)}
  \crefmultiformat{equation}{(#2#1#3}{,~#2#1#3)}{,~#2#1#3}{,~#2#1#3)}
  \crefrangemultiformat{equation}{(#3#1#4--#5#2#6}{,~#3#1#4--#5#2#6)}{,~#3#1#4--#5#2#6}{,~#3#1#4--#5#2#6)}
  
  \crefname{figure}{Figure}{Figures}
  \crefrangeformat{figure}{Figures~#3#1#4--#5#2#6}
  \crefmultiformat{figure}{Figures~#2#1#3}{ and~#2#1#3}{, #2#1#3}{ and~#2#1#3}
  \crefrangemultiformat{figure}{Figures~#3#1#4--#5#2#6}{ and~#3#1#4--#5#2#6}{, #3#1#4--#5#2#6}{ and~#3#1#4--#5#2#6}
  
  \crefname{table}{Table}{Tables}
  \crefrangeformat{table}{Tables~#3#1#4--#5#2#6}
  \crefmultiformat{table}{Tables~#2#1#3}{ and~#2#1#3}{, #2#1#3}{ and~#2#1#3}
  \crefrangemultiformat{table}{Tables~#3#1#4--#5#2#6}{ and~#3#1#4--#5#2#6}{, #3#1#4--#5#2#6}{ and~#3#1#4--#5#2#6}
  
  \crefname{algorithm}{Algorithm}{Algorithms}
  \crefrangeformat{algorithm}{Algorithms~#3#1#4--#5#2#6}
  \crefmultiformat{algorithm}{Algorithms~#2#1#3}{ and~#2#1#3}{, #2#1#3}{ and~#2#1#3}
  \crefrangemultiformat{algorithm}{Algorithms~#3#1#4--#5#2#6}{ and~#3#1#4--#5#2#6}{, #3#1#4--#5#2#6}{ and~#3#1#4--#5#2#6}
  
  %\newtheorem{theorem}{Theorem}
  %\newtheorem{lemma}[theorem]{Lemma}
  \ifcasscclass
    % «cas-sc» fornisce \newdefinition
    \newdefinition{dfn}{Definition}
    \newdefinition{rmk}{Remark}
    \newdefinition{hyp}{Hypotesis}
  \else
    % Con «article» (Articolo) si usano gli ambienti del kernel, con
    % contatore condiviso e «Assumption» come titolo di «hyp»
    \newtheorem{dfn}{Definition}
    \newtheorem{rmk}[dfn]{Remark}
    \newtheorem{hyp}[dfn]{Assumption}
  \fi
  %\newproof{pf}{Proof}
  %\newproof{pot}{Proof of Theorem \ref{thm}}
  
  \crefname{dfn}{Definition}{Definitions}
  \crefrangeformat{dfn}{Definitions~#3#1#4--#5#2#6}
  \crefmultiformat{dfn}{Definitions~#2#1#3}{ and~#2#1#3}{, #2#1#3}{ and~#2#1#3}
  \crefrangemultiformat{dfn}{Definitions~#3#1#4--#5#2#6}{ and~#3#1#4--#5#2#6}{, #3#1#4--#5#2#6}{ and~#3#1#4--#5#2#6}
  \crefname{hyp}{Hypothesis}{Hypotheses}
  
  \crefrangeformat{hyp}{Hypotheses~#3#1#4--#5#2#6}
  \crefmultiformat{hyp}{Hypotheses~#2#1#3}{ and~#2#1#3}{, #2#1#3}{ and~#2#1#3}
  \crefrangemultiformat{hyp}{Hypotheses~#3#1#4--#5#2#6}{ and~#3#1#4--#5#2#6}{, #3#1#4--#5#2#6}{ and~#3#1#4--#5#2#6}
  
  \crefname{rmk}{Remark}{Remarks}
  \crefrangeformat{rmk}{Remarks~#3#1#4--#5#2#6}
  \crefmultiformat{rmk}{Remarks~#2#1#3}{ and~#2#1#3}{, #2#1#3}{ and~#2#1#3}
  \crefrangemultiformat{rmk}{Remarks~#3#1#4--#5#2#6}{ and~#3#1#4--#5#2#6}{, #3#1#4--#5#2#6}{ and~#3#1#4--#5#2#6}
  
  % \newcommand\tikzfigure[5][0em]{%
  %   \begin{figure}[%
  %     pos=#2,%
  %     width=#4,%
  %     align=\centering%
  %   ]%
  %     \centering%
  %     \vspace*{#1}%
  %     \IfSubStr{#2}{p}{%
  %       \leavevmode\stepcounter{figposcount}%
  %       \zsaveposx{figpos-\thefigposcount}%
  %       \makebox[0pt][l]{%
  %         \hspace*{-\dimexpr\zposx{figpos-\thefigposcount}sp\relax}%
  %         \hspace*{\dimexpr0.5\paperwidth\relax}%
  %         \makebox[0pt][c]{%
  %           \scalebox{0.86}{\includegraphics{figs/#5.pdf}}%
  %         }%
  %       }
  %     }{%
  %       \includegraphics{figs/#5.pdf}%
  %     }%
  %     \caption{%
  %       \hyphenpenalty=0%
  %       \exhyphenpenalty=0%
  %       #3%
  %     }%
  %     \protect\label{#5}%
  %   \end{figure}%
  % }
%\endgroup

%\begingroup Versione per la classe «cas-sc»
  % 
  % %\begingroup For image generation
  %   %\begingroup Tipografia
  %     \newcommand\textsize[1]{%
  %       \pgfmathparse{#1*1.2}%
  %       \fontsize{#1pt}{\pgfmathresult pt}\selectfont%
  %     }
  % 
  %     \DeclareMathOperator\FRC{\mathcal R}
  % 
  %     \newif\ifcasscclass
  %     \makeatletter
  %     \@ifclassloaded{cas-sc}{\casscclasstrue}{\casscclassfalse}
  %     \makeatother
  % 
  %     \providecommand\Ifthispageodd[2]{\checkoddpage\ifoddpage#1\else#2\fi}
  %   %\endgroup
  % 
  %   %\begingroup TikZ
  %     %\begingroup Didascalie
  %       \newcommand\newcaption[2]{%
  %         \captionsetup{width=#2} % aboveskip=-.2cm,
  %         \caption{#1}%
  %       }
  %       
  %       \newcommand\newsubcaption[2][ ]{%
  %         \refstepcounter{subfigure}%
  %         (\thesubfigure)#1#2%
  %       }
  %     %\endgroup
  %     
  %     %\begingroup RedefineTikZbounds
  %       % La differenza è che colla prima bisogna catturare anche le etichette perché si crea un pdf isolato mentre colla non è importante perché viene tutto compilato sul momento
  % 
  %       %\begingroup Versione valida con esternalizzazione [ma anche senza]
  %         \newcommand\redefineTikZbounds[2]{
  %           \coordinate (RealSW) at (current bounding box.south west);
  %           \coordinate (RealNE) at (current bounding box.north east);
  %           \coordinate (Center) at ($(#1.west)!0.5!(#2.east)$);
  %     
  %           \pgfresetboundingbox 
  % 
  %           \useasboundingbox 
  %             let \p1=(RealSW),\p2=(RealNE),\p3=(Center),
  %             \n1={max(\x3-\x1,\x2-\x3)} in
  %             (\x3-\n1,\y1) rectangle (\x3+\n1,\y2+.25em);
  % 
  %           % \path[use as bounding box] 
  %           %   let \p1=(RealSW),\p2=(RealNE),\p3=(Center),
  %           %   \n1={max(\x3-\x1,\x2-\x3)} in
  %           %   (\x3-\n1,\y1) rectangle (\x3+\n1,\y2+.25em);
  %           % Lo «.25em» discende dalla traslazione imposta all'etichetta della scala verticale mediante «yScaleLabelShift»
  %         } % Scatola delimitante simmetrica
  %       %\endgroup
  % 
  %       %\begingroup Versione valida senza esternalizzazione
  %         % \newcommand\redefineTikZbounds[2]{
  %         %   % Salva le coordinate della cima e del fondo della figura corrente (cosí da mantenere i titoli e le didascalie)
  %         %   \coordinate (Cima) at (current bounding box.north);
  %         %   \coordinate (Fondo) at (current bounding box.south);
  %     
  %         %   \pgfresetboundingbox % Resetta la scatola dei limiti
  %     
  %         %   \useasboundingbox % Definisce la nuova scatola come un rettangolo
  %         %     (#1.west |- Cima)  % con ascissa sinistra data da «Col1 c1r1.west», 
  %         %     rectangle      % ascissa destra data da «Col2 c1r1.east»
  %         %     (#2.east |- Fondo); % ordinate date dalla cima e fondo originali
  %         %     % L'operatore «|-» in «(A |- B)» prende come ascissa quella del punto A mentre come coordinata quella del punto B
  %         % }
  %       %\endgroup
  %     %\endgroup
  % 
  %     %\begingroup TikZBox
  %       % Parzialmente ispirato da https://gist.github.com/Caramdir/899107
  % 
  %       \makeatletter
  %       \newsavebox\tikzBufferBox
  %       \xdef\offsetX{0pt}
  %       \xdef\offsetY{0pt}
  %       
  %       % \tikzbox[#1]{#2}{#3}[#4]{#5}
  %       % #1 -> Ancora rispetto cui posizionare l'intera scatola (di base è «south west», ossia l'angolo in basso a sinistra)
  %       % #2 -> Nodo dove posizionare l'origine (per la natura di come LaTeX riceve gli argomenti è buon pratica passare questo circondato da parentesi graffe)
  %       % #3 -> Nome della scatola limitante definita dall'ambiente «scope»
  %       % #4 -> Argomenti opzionali da passare all'ambiente «scope»; di base sono vuoti
  %       % #5 -> Contenuto della scatola circondato dall'ambiente «scope»
  %       % Es. \tikzbox{{(0,0)}}{figuraTikZ}{blocco TikZ}
  % 
  %       % Macro di apertura che controlla se v'è l'argomento facoltativo dell'ancora
  %       \def\tikzbox{\pgfutil@ifnextchar[{\tikzbox@anchor}{\tikzbox@anchor[south west]}}
  %       % La funzione «\pgfutil@ifnextchar» serve per controllare se il prossimo carattere è l'apertura di una paretensi graffa «[»
  %       
  %       % Macro intermedia per cogliere gli argomenti obbligatori e controllare quelli successivi facoltativi
  %       \def\tikzbox@anchor[#1]#2#3{\pgfutil@ifnextchar[{\tikzbox@opt[#1]#2#3}{\tikzbox@opt[#1]#2#3[]}}
  % 
  %       % Macro degli argomenti obbligatori
  %       \long\def\tikzbox@opt[#1]#2#3[#4]#5{%
  %         \node[ % Misurazione della scatola in un nodo invisibile
  %           overlay,%
  %           opacity=0,%
  %           inner sep=0pt,%
  %           outer sep=0pt%
  %         ] {%
  %           \sbox{\tikzBufferBox}{%
  %             \begin{tikzpicture}%
  %               \begin{scope}[local bounding box=measuringTikZBox]%
  %                 #5%
  %               \end{scope}%
  %               \pgfpointanchor{measuringTikZBox}{#1}%
  %               \pgfgetlastxy{\tempX}{\tempY}% 
  %               \xdef\offsetX{\tempX}%
  %               \xdef\offsetY{\tempY}%
  %             \end{tikzpicture}%
  %           }%
  %         };%
  %         %
  %         % Disegno reale traslato
  %         \begin{scope}[%
  %           local bounding box=#3,%
  %           shift={#2},%
  %           xshift=-\offsetX,%
  %           yshift=-\offsetY,%
  %           #4%
  %         ]%
  %           #5%
  %         \end{scope}%
  %       }%
  %       \makeatother
  %     %\endgroup
  % 
  %     %\begingroup Chiavi-Valori
  %       %\begingroup Testo
  %         % La dimensione del testo di base è 11pt
  %         \newcommand\titleSize{11}
  %         \newcommand\subCaptionSize{10}
  %         \newcommand\ParameterSize{9}
  %         \newcommand\labelSize{10}
  %         \newcommand\tickSize{9}
  %         \newcommand\legendSize{9}
  %       %\endgroup
  %     
  %       %\begingroup Infopagine
  %         %\begingroup Chiavi condivise
  %         \newcounter{lpc} % «loopcounter»
  %         \newcommand\infopageSharedKeys{
  %           %\begingroup Testo
  %             \renewcommand\titleSize{10}
  %             \renewcommand\subCaptionSize8
  %             \renewcommand\ParameterSize7
  %             \renewcommand\labelSize9
  %             \renewcommand\tickSize6
  %             \renewcommand\legendSize7
  %           %\endgroup
  %     
  %           %\begingroup Dimensioni
  %             \pgfmathsetmacro\plotWidth7
  %             \pgfmathsetmacro\plotHeight{.618*\plotWidth}
  %             \pgfmathsetmacro\halfPlotHeight{\plotHeight/2}
  %             %
  %             \pgfmathparse{0.825*\plotWidth}
  %             \pgfmathsetmacro\plotWidth\pgfmathresult
  %             \pgfmathsetmacro\halfWidth{\plotWidth/2}
  %             \pgfmathsetmacro\plotLastWidth{1.25*\plotWidth}
  %     
  %             \pgfmathsetmacro\xPlotSep{1}
  %             \pgfmathsetmacro\yPlotSep{1.25}
  %             \pgfmathsetmacro\yGroupSep{2.75}
  %           %\endgroup
  %     
  %           %\begingroup Legende
  %             \newcommand\placeLegend[5]{%
  %               /pgfplots/plotLegend,
  %               at={([shift={(##4,##5)}]axis description cs:##1,##2)},
  %               anchor=##3,
  %             }
  %             %
  %             \pgfplotsset{%
  %               plotLegend/.style={%
  %                 font=\textsize{\legendSize},
  %                 /pgfplots/legend cell align=left,
  %                 % align=left,
  %                 % draw=white!15!black,
  %                 draw=none, % Bordo invisibile
  %                 fill=white,
  %                 fill opacity=0.5,
  %                 text opacity=1,
  %                 % text width=2cm,
  %               },
  %               plotLegendNW/.style={legend style={\placeLegend{0}{1}{north west}{0em}{0em}}},
  %               plotLegendNE/.style={legend style={\placeLegend{1}{1}{north east}{0em}{0em}}},
  %               plotLegendSW/.style={legend style={\placeLegend{0}{0}{south west}{0em}{0em}}},
  %               plotLegendSE/.style={legend style={\placeLegend{1}{0}{south east}{0em}{0em}}},
  %               %
  %               % plotLegendW/.style={legend style={plotLegend,at={(0,.5)},anchor=west}},
  %               %
  %               % plotLegendNW/.style={legend style={plotLegend,legend pos=north west}},
  %               % plotLegendNE/.style={legend style={plotLegend,legend pos=north east}},
  %               % plotLegendSW/.style={legend style={plotLegend,legend pos=south west}},
  %               % plotLegendSE/.style={legend style={plotLegend,legend pos=south east}},
  %               %
  %               % plotLegendNWcloser/.style={legend style={plotLegend,at={(0,1)},anchor=north west}},
  %               % plotLegendNEcloser/.style={legend style={plotLegend,at={(1,1)},anchor=north east}},
  %               % plotLegendSWcloser/.style={legend style={plotLegend,at={(0,0)},anchor=south west}},
  %               % plotLegendSEcloser/.style={legend style={plotLegend,at={(0,1)},anchor=south east}},
  %             }
  %             %
  %             \pgfmathsetmacro\lineLegendLength{1}
  %             \pgfmathsetmacro\areaLegendHeigth{.25}
  %             \pgfplotsset{
  %               legend image code/.code={
  %                 \draw [mark repeat=2,mark phase=2,##1]
  %                   plot coordinates {
  %                     (0em,0em)
  %                     (.5*\lineLegendLength em,0em)
  %                     (\lineLegendLength em,0em)
  %                   };
  %               },
  %               /pgfplots/area legend/.style={
  %                 legend image code/.code={
  %                   \draw [##1,fill,draw=none] 
  %                     (0em,-\areaLegendHeigth em) 
  %                     rectangle 
  %                     (\lineLegendLength em,\areaLegendHeigth em);
  %                 },
  %               },
  %             }
  %             %
  %             \newcommand\ExactHistogramLegend{Exact histogram}
  %             \newcommand\ApproxHistogramLegend{Approx. histogram}
  %             \newcommand\RealHistogramLegend{Real histogram}
  %             %
  %             \newcommand\ExactBLNFitLegend{Exact BLN fit}
  %             \newcommand\ApproxBLNFitLegend{Approx BLN fit}
  %             \newcommand\RealBLNFitLegend{Real BLN fit}
  %             %
  %             \newcommand\ExactEmpiricCCDFLegend{Exact empiric CCDF}
  %             \newcommand\ApproxEmpiricCCDFLegend{Approx. empiric CCDF}
  %             \newcommand\RealEmpiricCCDFLegend{Real empiric CCDF}
  %             %
  %             \newcommand\ExactParetoFitLegend{Exact Pareto fit}
  %             \newcommand\ApproxParetoFitLegend{Approx. Pareto fit}
  %             \newcommand\RealParetoFitLegend{Real Pareto fit}
  %             %
  %             \newcommand\ExactScatterLegend{Exact scatter}
  %             \newcommand\ApproxScatterLegend{Approx. scatter}
  %             \newcommand\RealScatterLegend{Real scatter}
  %             %
  %             \newcommand\ClasskILegend{Class k=283}
  %             \newcommand\ClasskIILegend{Class k=83}
  %             \newcommand\ClasskIIILegend{Class k=43}
  %             \newcommand\ClasskIVLegend{Class k=10}
  %           %\endgroup
  %     
  %           %\begingroup Etichette vuote
  %             \pgfplotsset{
  %               yLabelEmpty/.style={ylabel={}},
  %               yTickLabelsEmpty/.style={yticklabels={}},
  %               yTickScaleLabelEmpty/.style={ytick scale label code/.code={}},
  %               xTickScaleLabelEmpty/.style={xtick scale label code/.code={}},
  %               emptyLeged/.style={
  %                 legend style={draw=none,fill=none,nodes={draw=none}},
  %                 legend image post style={opacity=0},
  %               },
  %             }
  %           %\endgroup
  %     
  %           %\begingroup Estetica
  %             \newcommand\smallminus{\scalebox{0.5}[1]{$-$}\hspace{-.11em}}
  %             \pgfplotsset{
  %               /pgfplots/log number format basis/.code 2 args={%
  %                 \pgfmathsetmacro\exponent{##2}%
  %                 \pgfmathparse{sign(\exponent)}%
  %                 \ifnum\pgfmathresult<0%
  %                   \pgfmathparse{abs(\exponent)}%
  %                   $##1^{\smash{\smallminus\pgfmathprintnumber{\pgfmathresult}}}$%
  %                 \else%
  %                   $##1^{\smash{\pgfmathprintnumber{\exponent}}}$%
  %                 \fi%
  %               },
  %               ytick scale label code/.code={%
  %                 \pgfmathsetmacro\exponent{##1}%
  %                 \pgfmathparse{sign(\exponent)}%
  %                 \ifnum\pgfmathresult<0
  %                   \pgfmathparse{abs(\exponent)}%
  %                   $\cdot 10^{\smash{\smallminus\pgfmathprintnumber{\pgfmathresult}}}$%
  %                 \else
  %                   $\cdot 10^{\smash{\pgfmathprintnumber{\exponent}}}$%
  %                 \fi
  %               },
  %             }
  %           %\endgroup
  %     
  %           %\begingroup Colori
  %             \definecolor{blueish}{rgb}{0.00000,0.44700,0.74100}%
  %             \definecolor{reddish}{rgb}{0.83529,0.36863,0.00000}%
  %           %\endgroup
  %     
  %           %\begingroup Commenti
  %             %\begingroup Impostazioni
  %             \tikzset{
  %               title/.style={
  %                 align=center,
  %                 anchor=south,
  %                 font=\bfseries\textsize{\titleSize},
  %               },
  %               fittedParameters/.style={
  %                 align=center,
  %                 font=\textsize{\ParameterSize},
  %               },
  %               varyingParameter/.style={
  %                 align=center,
  %                 % rotate=90,
  %                 font=\textsize{\subCaptionSize},
  %               },
  %               subCaption/.style={
  %                 % text width=\plotWidth cm,
  %                 align=center,
  %                 font=\textsize{\subCaptionSize},
  %               },
  %             }
  %             \newcommand\codetocolor[1]{\ifcase##1blueish\or reddish\else black\fi}
  %             %\endgroup
  %     
  %             %\begingroup Parametri Sardegna
  %             \newcommand\betaReal{\num{1.273}}
  %             \newcommand\xiReal{\num{0.929}}
  %             \newcommand\mIReal{\num{3.227}}
  %             \newcommand\sIReal{\num{0.408}}
  %             \newcommand\mIIReal{\num{4.099}}
  %             \newcommand\sIIReal{\num{0.556}}
  %             %\endgroup
  %     
  %             %\begingroup Parametri Pareto
  %             \newcommand\codetoParetovalue[1]{
  %               \ifcase##1\betaExact\or\betaAprx\else\betaReal\fi
  %             }
  %             \newcommand\writeParetoParameters[7][center]{
  %               \node[fittedParameters,anchor=##1] at ([shift={(##5,##6)}]##4) {%
  %                 \(\begin{aligned}
  %                   \color{\codetocolor{##2}}\beta
  %                   &\color{\codetocolor{##2}}=\codetoParetovalue{##2}\\[##7]
  %                   \color{\codetocolor{##3}}\beta
  %                   &\color{\codetocolor{##3}}=\codetoParetovalue{##3}
  %                 \end{aligned}\)
  %               };
  %             }
  %             %\endgroup
  % 
  %             %\begingroup % Vecchia versione in cui si scrivevano esplicitamente tutte le 9 combinazioni possibili
  %             % \newcommand\writeParameters[6]{%
  %             %   \node[parameter] at ([shift={(##4,##5)}]##1) {%
  %             %     \ifcase##2% Caso 0
  %             %       \ifcase##3% Caso 0
  %             %         \(\begin{aligned}
  %             %           \color{blueish}\beta&\color{blueish}=\betaExact\\[##6]
  %             %           \color{blueish}\beta&\color{blueish}=\betaExact
  %             %         \end{aligned}\)
  %             %       \or% Caso 1
  %             %         \(\begin{aligned}
  %             %           \color{blueish}\beta&\color{blueish}=\betaExact\\[##6]
  %             %           \color{reddish}\beta&\color{reddish}=\betaAprx% Caso 1
  %             %         \end{aligned}\)
  %             %       \else% Caso 2
  %             %         \(\begin{aligned}
  %             %           \color{blueish}\beta&\color{blueish}=\betaExact\\[##6]
  %             %           \color{black}\beta&\color{black}=\betaReal% Caso 2
  %             %         \end{aligned}\)
  %             %       \fi
  %             %     \or% Caso 1
  %             %       \ifcase##3% Caso 0
  %             %         \(\begin{aligned}
  %             %           \color{reddish}\beta&\color{reddish}=\betaAprx\\[##6]
  %             %           \color{blueish}\beta&\color{blueish}=\betaExact
  %             %         \end{aligned}\)
  %             %       \or% Caso 1
  %             %         \(\begin{aligned}
  %             %           \color{reddish}\beta&\color{reddish}=\betaAprx\\[##6]
  %             %           \color{reddish}\beta&\color{reddish}=\betaAprx
  %             %         \end{aligned}\)
  %             %       \else% Caso 2
  %             %         \(\begin{aligned}
  %             %           \color{reddish}\beta&\color{reddish}=\betaAprx\\[##6]
  %             %           \color{black}\beta&\color{black}=\betaReal
  %             %         \end{aligned}\)
  %             %       \fi
  %             %     \else% Caso 2
  %             %       \ifcase##3% Caso 0
  %             %         \(\begin{aligned}
  %             %           \color{black}\beta&\color{black}=\betaReal\\[##6]
  %             %           \color{blueish}\beta&\color{blueish}=\betaExact
  %             %         \end{aligned}\)
  %             %       \or% Caso 1
  %             %         \(\begin{aligned}
  %             %           \color{black}\beta&\color{black}=\betaReal\\[##6]
  %             %           \color{reddish}\beta&\color{reddish}=\betaAprx
  %             %         \end{aligned}\)
  %             %       \else% Caso 2
  %             %         \(\begin{aligned}
  %             %           \color{black}\beta&\color{black}=\betaReal\\[##6]
  %             %           \color{black}\beta&\color{black}=\betaReal
  %             %         \end{aligned}\)
  %             %       \fi
  %             %     \fi
  %             %   };
  %             % }
  %             %\endgroup
  %     
  %             %\begingroup Parametri bilognormale
  %             \newcommand\codetoparameter[1]{
  %               \ifcase##1\xi  % Caso 0
  %               \or\mu_1     % Caso 1
  %               \or\sigma_1   % Caso 2
  %               \or\mu_2     % Caso 3
  %               \else\sigma_2\fi % Caso 4
  %             }
  %             \newcommand\codetoBiLognormalvalue[2]{
  %               \ifcase##1 % Caso 0
  %                 \ifcase##2\xiExact % Caso 0
  %                 \or\mIExact    % Caso 1
  %                 \or\sIExact    % Caso 2
  %                 \or\mIIExact    % Caso 3
  %                 \else\sIIExact\fi % Caso 4
  %               \or % Caso 1
  %                 \ifcase##2\xiAprx % Caso 0
  %                 \or\mIAprx     % Caso 1
  %                 \or\sIAprx     % Caso 2
  %                 \or\mIIAprx    % Caso 3
  %                 \else\sIIAprx\fi  % Caso 4
  %               \else % Caso 2
  %                 \ifcase##2\xiReal % Caso 0
  %                 \or\mIReal     % Caso 1
  %                 \or\sIReal     % Caso 2
  %                 \or\mIIReal    % Caso 3
  %                 \else\sIIReal\fi  % Caso 4
  %               \fi
  %             }
  %             \newcommand\matrixline[2]{\codetoparameter{##2}&=\codetoBiLognormalvalue{##1}{##2}}
  %             \def\lineloop##1##2{%
  %               \matrixline{##1}{\number\value{lpc}}\stepcounter{lpc}%
  %               \ifnum\number\value{lpc}<5\\[##2]\lineloop{##1}{##2}\fi%
  %             } % https://tex.stackexchange.com/questions/491358/tikz-foreach-to-work-with-the-aligned-environment
  %             \newcommand\writeBiLognormalParameters[6][center]{
  %               \node[fittedParameters,anchor=##1] at ([shift={(##4,##5)}]##3) {%
  %                 \setcounter{lpc}{0}
  %                 \(\color{\codetocolor{##2}}
  %                 \begin{aligned}
  %                   \lineloop{##2}{##6}
  %                 \end{aligned}\)
  %               };
  %             }
  %             %\endgroup
  %     
  %             %\begingroup Altre etichette
  %             \newcommand\writeSizeEvolutionDescriptions[1][5]{%
  %               \foreach \col/\clr/\txt in {%
  %                 1/blueish/Esatto,
  %                 2/reddish/Approssimato
  %               }{
  %                 \node[
  %                   fittedParameters,
  %                   anchor=south,
  %                   yshift=.1em,
  %                 ] at (Row##1 c\col r1.south) {\color{\clr}\txt};
  %               }
  %             }
  %     
  %             \newcommand\writeVaryingParameter[5]{%
  %               \foreach \r in {1,...,5}{
  %                 \node[
  %                   varyingParameter,
  %                   anchor=##3,
  %                 ] at ([shift={(##4,##5)}]Row\r\space c##1r1.##2)
  %                 {\(\varyingParameterSymbol=\varyingParameterList\r\)};
  %               }
  %             }
  %             %\endgroup
  %           %\endgroup
  %     
  %           %\begingroup Sottodidascalie
  %             \newcommand\writeSubCaption[5]{%
  %               \node[subCaption,anchor=##3]
  %               at ([shift={(##4,##5)}]##1.##2)
  %               {\newsubcaption[]{}};
  %             }
  %             \newcommand\writeSubCaptions[5][3]{%
  %               \foreach \r in {1,...,5}{
  %                 \foreach \c in {1,...,##1}{
  %                   \writeSubCaption{Row\r\space c\c r1}{##2}{##3}{##4}{##5}
  %                 }
  %               }
  %             }
  %           %\endgroup
  %         }
  %         %\endgroup
  %     
  %         %\begingroup Chiavi del primo tipo d'infopagina
  %         \newcommand\infopageKeysI{
  %           \infopageSharedKeys
  %     
  %           %\begingroup Gruppi
  %             \ifpageisodd
  %               \pgfplotsset{/pgfplots/group/yDescPos/.style={y descriptions at=edge right}}
  %             \else
  %               \pgfplotsset{/pgfplots/group/yDescPos/.style={y descriptions at=edge left}}
  %             \fi
  %             \pgfplotsset{
  %               /pgfplots/group/plotGroup/.style={
  %                 horizontal sep=\xPlotSep em,
  %                 vertical sep=\yPlotSep em,
  %                 yDescPos,
  %                 group size=3 by 1,
  %               },
  %               /pgfplots/group/plotLastGroup/.style={
  %                 plotGroup,
  %                 x descriptions at=edge bottom,
  %               },
  %             }
  %           %\endgroup
  %     
  %           %\begingroup Stili
  %             %\begingroup Selezione «recto»/«verso»
  %             \ifpageisodd
  %               \pgfplotsset{%
  %                 yLineAxis/.style={axis y line*=right},
  %                 /tikz/yLabelShift/.style={%
  %                   yshift=-1.25em,
  %                   at={(axis description cs:1,0.5)},
  %                   anchor=north,
  %                   % rotate=180,
  %                 },
  %                 /tikz/yScaleLabelShift/.style={%
  %                   yshift=.25em,
  %                   at={(1,1)},
  %                   anchor=south east,
  %                 },
  %                 /tikz/yTickLabelShift/.style={xshift=-0.15em},
  %               }
  %             \else
  %               \pgfplotsset{%
  %                 yLineAxis/.style={axis y line*=left},
  %                 /tikz/yLabelShift/.style={%
  %                   yshift=1.25em,
  %                   at={(axis description cs:0,0.5)},
  %                   anchor=south,
  %                   % rotate=180,
  %                 },
  %                 /tikz/yScaleLabelShift/.style={%
  %                   yshift=.25em,
  %                   at={(0,1)},
  %                   anchor=south west,
  %                 },
  %                 /tikz/yTickLabelShift/.style={xshift=+0.15em},
  %               }
  %             \fi
  %             %\endgroup
  %             %
  %             %\begingroup Grafici
  %             \pgfplotsset{
  %               plotRow/.style={
  %                 scale only axis,
  %                 width=\plotWidth cm,height=\plotHeight cm,
  %                 xlabel style={font=\color{white!15!black}\textsize{\labelSize}},
  %                 ylabel style={font=\color{white!15!black}\textsize{\labelSize},yLabelShift},
  %                 xticklabel style={font=\textsize{\tickSize}},
  %                 yticklabel style={font=\textsize{\tickSize},yTickLabelShift},
  %                 y tick scale label style={yScaleLabelShift},
  %                 % axis x line*=bottom,yLineAxis,
  %                 enlarge x limits=0.1,enlarge y limits=0.1,
  %                 xmajorgrids,ymajorgrids,
  %                 xminorticks=true,yminorticks=true,
  %                 scale ticks above exponent=2,
  %                 axis background/.style={fill=white},
  %                 grid style={dashed}
  %               },
  %               % xlabel123/.style={xlabel={\(h\)}},
  %               xlabel3/.style={xlabel={\(h\)}},
  %               ylabel23/.style={ylabel={\(\FRC(h)\)}},
  %               row1Keys/.style={
  %                 % ylabel={\(\underhat{\bar l}_N(h)\)},
  %                 ylabel={\(\bar l_N(h)\)},
  %                 % enlarge x limits={lower,value=.05},
  %                 enlarge x limits=.05,
  %                 enlarge y limits={upper,value=.05},
  %                 xmode=linear,
  %                 log origin=infty,
  %               },
  %               row2KeysWithLabels/.style={
  %                 ylabel23,
  %                 xmode=linear,ymode=log
  %               },
  %               row2Keys/.style={
  %                 ylabel23,
  %                 xmode=linear,ymode=log,
  %                 xticklabels={},
  %               },
  %               row3Keys/.style={
  %                 xlabel3,ylabel23,
  %                 xmode=linear,ymode=log
  %               },
  %               row4Keys/.style={
  %                 xlabel={\(\xi\)},ylabel={\(s(\xi)\)},
  %                 xmode=log,ymode=log
  %               },
  %               row5KeysLin/.style={
  %                 width=\plotLastWidth cm,
  %                 xlabel={\(t\)},ylabel={\(s(t)\)},
  %                 enlarge x limits=0,
  %                 enlarge y limits=0.05,
  %                 x tick scale label style={
  %                   yshift=-1em,
  %                   at={(.95,0)},
  %                   anchor=north,
  %                 },
  %               },
  %               row5Keys/.style={
  %                 row5KeysLin,
  %                 ymode=log,
  %               },
  %             }
  %             %\endgroup
  %           %\endgroup
  %         }
  %         %\endgroup
  %     
  %         %\begingroup Chiavi del secondo tipo d'infopagina
  %         \newcommand\infopageKeysII{
  %           \infopageSharedKeys
  %     
  %           %\begingroup Gruppi
  %             \ifpageisodd
  %               \pgfplotsset{/pgfplots/group/yDescPos/.style={y descriptions at=edge right}}
  %             \else
  %               \pgfplotsset{/pgfplots/group/yDescPos/.style={y descriptions at=edge left}}
  %             \fi
  %             \pgfplotsset{
  %               /pgfplots/group/plotGroup/.style={
  %                 horizontal sep=\xPlotSep em,
  %                 vertical sep=\yPlotSep em,
  %                 y descriptions at=all,
  %                 group size=3 by 1,
  %               },
  %               /pgfplots/group/plotLastGroup/.style={
  %                 plotGroup,
  %                 x descriptions at=edge bottom,
  %               },
  %             }
  %           %\endgroup
  %     
  %           %\begingroup Stili
  %             %\begingroup Selezione «recto»/«verso»
  %             \ifpageisodd
  %               \pgfplotsset{
  %                 yLabel/.style={
  %                   ylabel style={
  %                     font=\color{white!15!black}\textsize{\labelSize},
  %                     yshift=-1.25em,
  %                     at={(axis description cs:1,0.5)},
  %                     anchor=north,
  %                     % rotate=180,
  %                   },
  %                 },
  %                 yLabelInv/.style={yLabelEmpty},
  %                 yScaleLabelShift/.style={
  %                   y tick scale label style={
  %                     yshift=.25em,
  %                     % xshift=1em,
  %                     at={(1,1)},
  %                     anchor=south east,
  %                   },
  %                 },
  %                 yScaleLabelShiftInv/.style={
  %                   y tick scale label style={
  %                     yshift=.25em,
  %                     % xshift=-1em,
  %                     at={(0,1)},
  %                     anchor=south west,
  %                   },
  %                 },
  %                 % yLineAxis/.style={axis y line*=right},
  %               }
  %             \else
  %               \pgfplotsset{
  %                 yLabel/.style={yLabelEmpty},
  %                 yLabelInv/.style={
  %                   ylabel style={
  %                     font=\color{white!15!black}\textsize{\labelSize},
  %                     yshift=1.25em,
  %                     at={(axis description cs:0,0.5)},
  %                     anchor=south,
  %                     % rotate=180,
  %                   },
  %                 },
  %                 yScaleLabelShift/.style={
  %                   y tick scale label style={
  %                     yshift=.25em,
  %                     % xshift=1em,
  %                     at={(1,1)},
  %                     anchor=south east,
  %                   },
  %                 },
  %                 yScaleLabelShiftInv/.style={
  %                   y tick scale label style={
  %                     yshift=.25em,
  %                     % xshift=-1em,
  %                     at={(0,1)},
  %                     anchor=south west,
  %                   },
  %                 },
  %                 % yLineAxis/.style={axis y line*=left},
  %               }
  %             \fi
  %             %\endgroup
  %             %
  %             %\begingroup Grafici
  %             \pgfplotsset{
  %               plotRow/.style={
  %                 scale only axis,
  %                 width=\plotWidth cm,height=\plotHeight cm,
  %                 xlabel style={font=\color{white!15!black}\textsize{\labelSize}},
  %                 % ylabel style={font=\color{white!15!black}\textsize{\labelSize},yLabelShift},
  %                 xticklabel style={font=\textsize{\tickSize}},
  %                 % axis x line*=bottom,yLineAxis,
  %                 enlarge x limits=0.1,enlarge y limits=0.1,
  %                 xmajorgrids,ymajorgrids,
  %                 xminorticks=true,yminorticks=true,
  %                 scale ticks above exponent=2,
  %                 axis background/.style={fill=white},
  %                 grid style={dashed},
  %               },
  %               plotLeftCol/.style={
  %                 yLabelInv,
  %                 yticklabel style={font=\textsize{\tickSize},xshift=0.15em},
  %                 yticklabel pos=left,
  %                 yScaleLabelShiftInv,
  %               },
  %               plotMiddleCol/.style={
  %                 yLabelEmpty,
  %                 yTickLabelsEmpty,
  %                 yTickScaleLabelEmpty,
  %               },
  %               plotRightCol/.style={
  %                 yLabel,
  %                 yticklabel style={font=\textsize{\tickSize},xshift=-0.15em},
  %                 yticklabel pos=right,
  %                 yScaleLabelShift,
  %               },
  %               xlabel3/.style={xlabel={\(h\)}},
  %               ylabel23/.style={ylabel={\(\FRC(h)\)}},
  %               row1Keys/.style={
  %                 ylabel={\(\bar l_N(h)\)},
  %                 enlarge x limits=.05,
  %                 enlarge y limits={upper,value=.05},
  %                 xmode=linear,
  %                 log origin=infty,
  %               },
  %               row2Keys/.style={
  %                 ylabel23,
  %                 xmode=linear,ymode=log,
  %                 % xticklabels={},
  %               },
  %               row3Keys/.style={
  %                 xlabel3,ylabel23,
  %                 xmode=linear,ymode=log,
  %               },
  %               row4KeysLin/.style={
  %                 xlabel={\(\xi\)},ylabel={\(s(\xi)\)},
  %                 xmode=log,
  %               },
  %               row4Keys/.style={row4KeysLin,ymode=log},
  %               row5KeysLin/.style={
  %                 width=\plotLastWidth cm,
  %                 xlabel={\(t\)},ylabel={\(s(t)\)},
  %                 enlarge x limits=0,
  %                 enlarge y limits=0.05,
  %                 x tick scale label style={
  %                   yshift=-1em,
  %                   at={(.95,0)},
  %                   anchor=north,
  %                 },
  %               },
  %               row5Keys/.style={
  %                 row5KeysLin,
  %                 ymode=log,
  %               },
  %             }
  %             %\endgroup
  %           %\endgroup
  %         }
  %         %\endgroup
  %       %\endgroup
  %     %\endgroup
  % 
  %     %\begingroup Diagrammi dei nodi
  %     \newcommand\nodeDiagramSettings{
  %       %\begingroup Dimensioni
  %         \pgfmathsetmacro\xPlotSep{10}
  %         \pgfmathsetmacro\yPlotSep{4}
  %         \pgfmathsetmacro\plotScale{.8}
  % 
  %         \pgfmathsetmacro\titleSize{11}
  %         \pgfmathsetmacro\commentSize{10}
  % 
  %         \pgfmathsetmacro\nodeSize{20pt} % Diametro del nodo
  %         \pgfmathsetmacro\starSize{2.25*\nodeSize}
  %         \pgfmathsetmacro\numberSize{.53*\nodeSize}
  %       %\endgroup
  % 
  %       %\begingroup Colori
  %         % \definecolor{cRed}{HTML}{D68490}
  %         % \definecolor{cPurple}{HTML}{7E6EAC}
  %         % \definecolor{cTeal}{HTML}{87BDB1}
  %         % \definecolor{cOrange}{HTML}{F2D0A4}
  %         % \definecolor{cEdgeLight}{HTML}{EEE8D5} 
  %         % \definecolor{cEdgeDark}{HTML}{998888}
  %         % \definecolor{cStar}{HTML}{FDF6E3}
  %         %
  %         \definecolor{color1}{HTML}{800000} % Deep Red
  %         \definecolor{color2}{HTML}{4682B4} % Steel Blue
  %         \definecolor{color3}{HTML}{6B8E23} % Olive Green
  %         \definecolor{color4}{HTML}{708090} % Slate Gray
  %         \definecolor{color5}{HTML}{BC8F8F} % Rosy Brown
  %         \definecolor{color6}{HTML}{008B8B} % Dark Cyan
  %         \definecolor{color7}{HTML}{B8860B} % Dark Goldenrod
  %       %\endgroup
  % 
  %       %\begingroup Testo
  %         \tikzset{
  %           titleStyle/.style={
  %             font=\textsize{\titleSize},
  %             align=center,
  %             anchor=south,
  %           },
  %           commentStyle/.style={font=\textsize{\commentSize}},
  %           arrowCommentStyle/.style={
  %             midway,
  %             commentStyle,
  %             left,
  %             sloped,
  %             below
  %           },
  %           numberStyle/.style={font=\textsize{\numberSize}}
  %         }
  %       %\endgroup
  % 
  %       %\begingroup Frecce
  %         \tikzset{
  %           transformingArrow/.style={
  %             % =Latex,
  %             ->,
  %             >=Stealth,
  %             dashed,
  %             line width=.8pt,
  %             black!90,
  %             shorten >=5pt,
  %             shorten <=10pt
  %           },
  %           indirectArrow/.style={
  %             draw=black!95,
  %             line width=.5pt,
  %           },
  %           nodeArrowExact/.style={
  %             indirectArrow,
  %             % ->,
  %             % -{Stealth[scale=1.2]},
  %             -{Triangle[scale=.75]},
  %             % draw=cEdgeLight,
  %             % draw=red!50,
  %             % shorten >=2pt,
  %             % shorten <=2pt
  %           },
  %           nodeArrowApprox/.style={
  %             indirectArrow,
  %             -,
  %             dashed,
  %           },
  %           transparentNodeArrow/.style={
  %             opacity=.4,
  %           },
  %           arrowAxis/.style={
  %             ->,
  %             line width=.3pt,
  %             >=Latex,
  %           },
  %         }
  %       %\endgroup
  % 
  %       %\begingroup Nodi
  %         \tikzset{
  %           nodeStyle/.style={
  %             circle,
  %             minimum size=\nodeSize,
  %             % draw=black!80,
  %             draw=black,
  %             line width=.5pt,
  %             inner sep=0pt,
  %           },
  %           starStyle/.style={
  %             star,
  %             star points=5,
  %             star point ratio=2,
  %             % fill=cStar,
  %             fill=black!50,
  %             fill opacity=.25,
  %             minimum size=\starSize,
  %             draw=none
  %           },
  %         }
  %       %\endgroup
  %     }
  %     %\endgroup
  % 
  %     %\begingroup GNU-Plot
  %       % \IfFileExists{\sourceauxpath.aux}{
  %       %   \externaldocument{\sourceauxpath}
  %       % }{\relax}
  %       
  %       \pgfkeys{/pgf/plot/gnuplot call={cd /d "\sourceoutdir" && gnuplot}}
  %     %\endgroup
  %   %\endgroup
  % %\endgroup
  % 
  % %\begingroup For the paper
  %   \def\tsc#1{\csdef{#1}{\textsc{\lowercase{#1}}\xspace}}
  %   \tsc{WGM}
  %   \tsc{QE}
  %   
  %   \newcommand\etal{et al.}
  %   \newcommand\KTMAS{\hyperlink{acrKineticTheoryMultiAgentSystems}{KTMAS}}
  %   
  %   % \newcommand\MathSpacing[3]{
  %   %   \thinmuskip = #1mu\relax
  %   %   \medmuskip = #2mu\relax
  %   %   \thickmuskip = #3mu\relax
  %   % }
  %   % \MathSpacing123
  %   
  %   \newcommand\RPlus{\mathbb R_+}
  %   \newcommand\NPlus{\mathbb N_+}
  %   \DeclareMathAlphabet\pazocal{OMS}{zplm}{m}{n}
  %   
  %   \DeclareMathOperator\Lognormal{Lognormal}
  %   \DeclareMathOperator\BiLognormal{BiLognormal}
  %   \DeclareMathOperator\Bernoulli{Bernoulli}
  %   \DeclareMathOperator\Exponential{Exp}
  %   \DeclareMathOperator\Student{Student}
  %   \DeclareMathOperator\Variance{var}
  %   \DeclareMathOperator\IC{IC}
  %   \DeclareMathOperator\Histogram{\mathcal B}
  %   \DeclareMathOperator\CCDF{\mathcal R}
  %     
  %   \newcommand\ThreeQuarters{{3\hspace{-0.15em}/\hspace{-0.15em}4}}
  %   \newcommand\OneQuarter{{1\hspace{-0.15em}/\hspace{-0.15em}4}}
  %   
  %   \newcommand\de[1][]{\mathrm{d}#1}
  %   \newcommand\derS[3][d]{\frac{\mathrm{#1}#2}{\mathrm{#1}#3}}
  %   \newcommand\derP[2]{\frac{\partial#1}{\partial#2}}
  %   \newcommand\underhat[1]{\underaccent{\hat}{#1}}
  %   \renewcommand\undertilde[1]{\underaccent{\tilde}{#1}}
  %   
  %   \makeatletter
  %   \newcommand\eqtag{\@ifstar{\eqtag@star}{\eqtag@nostar}}
  %   
  %   \newcounter{eqtag}
  %   \newcommand\eqtag@nostar[2]{%
  %     \begin{equation*}%
  %       \refstepcounter{eqtag}%
  %       % \protected@edef\@currentlabeltype{eqtag}%
  %       \protected@edef\@currentlabel{#1}%
  %       #2%
  %       \tag*{#1}%
  %     \end{equation*}%
  %   }
  %   \newcommand\eqtag@star[2]{%
  %     \begin{equation*}%
  %       \refstepcounter{eqtag}%
  %       % \protected@edef\@currentlabeltype{eqtag}%
  %       \protected@edef\@currentlabel{#1}%
  %       \let\label\ltx@label%
  %       #2%
  %       % \notag%
  %     \end{equation*}%
  %   }
  %   \makeatother
  %   
  %   \newcommand\crefsubfigInt[4][-]{Figures~\ref{#2}#3#1\ref{#2}#4}
  %   \newcommand\crefsubfigSep[4]{Figures~\ref{#1}#2\foreach \l in#3{,~\ref{#1}\l}~and~\ref{#1}#4}
  %   
  %   \crefname{section}{§}{§§}
  %   \crefrangeformat{section}{§§~#3#1#4--#5#2#6}
  %   \crefmultiformat{section}{§§~#2#1#3}{ and~#2#1#3}{, #2#1#3}{ and~#2#1#3}
  %   \crefrangemultiformat{section}{§§~#3#1#4--#5#2#6}{ and~#3#1#4--#5#2#6}{, #3#1#4--#5#2#6}{ and~#3#1#4--#5#2#6}
  %   
  %   \crefname{equation}{}{}
  %   \crefrangeformat{equation}{(#3#1#4--#5#2#6)}
  %   \crefmultiformat{equation}{(#2#1#3}{,~#2#1#3)}{,~#2#1#3}{,~#2#1#3)}
  %   \crefrangemultiformat{equation}{(#3#1#4--#5#2#6}{,~#3#1#4--#5#2#6)}{,~#3#1#4--#5#2#6}{,~#3#1#4--#5#2#6)}
  %   
  %   \crefname{figure}{Figure}{Figures}
  %   \crefrangeformat{figure}{Figures~#3#1#4--#5#2#6}
  %   \crefmultiformat{figure}{Figures~#2#1#3}{ and~#2#1#3}{, #2#1#3}{ and~#2#1#3}
  %   \crefrangemultiformat{figure}{Figures~#3#1#4--#5#2#6}{ and~#3#1#4--#5#2#6}{, #3#1#4--#5#2#6}{ and~#3#1#4--#5#2#6}
  %   
  %   \crefname{table}{Table}{Tables}
  %   \crefrangeformat{table}{Tables~#3#1#4--#5#2#6}
  %   \crefmultiformat{table}{Tables~#2#1#3}{ and~#2#1#3}{, #2#1#3}{ and~#2#1#3}
  %   \crefrangemultiformat{table}{Tables~#3#1#4--#5#2#6}{ and~#3#1#4--#5#2#6}{, #3#1#4--#5#2#6}{ and~#3#1#4--#5#2#6}
  %   
  %   \crefname{algorithm}{Algorithm}{Algorithms}
  %   \crefrangeformat{algorithm}{Algorithms~#3#1#4--#5#2#6}
  %   \crefmultiformat{algorithm}{Algorithms~#2#1#3}{ and~#2#1#3}{, #2#1#3}{ and~#2#1#3}
  %   \crefrangemultiformat{algorithm}{Algorithms~#3#1#4--#5#2#6}{ and~#3#1#4--#5#2#6}{, #3#1#4--#5#2#6}{ and~#3#1#4--#5#2#6}
  %   
  %   %\newtheorem{theorem}{Theorem}
  %   %\newtheorem{lemma}[theorem]{Lemma}
  %   \newdefinition{dfn}{Definition}
  %   \newdefinition{rmk}{Remark}
  %   \newdefinition{hyp}{Hypotesis}
  %   %\newproof{pf}{Proof}
  %   %\newproof{pot}{Proof of Theorem \ref{thm}}
  %   
  %   \crefname{dfn}{Definition}{Definitions}
  %   \crefrangeformat{dfn}{Definitions~#3#1#4--#5#2#6}
  %   \crefmultiformat{dfn}{Definitions~#2#1#3}{ and~#2#1#3}{, #2#1#3}{ and~#2#1#3}
  %   \crefrangemultiformat{dfn}{Definitions~#3#1#4--#5#2#6}{ and~#3#1#4--#5#2#6}{, #3#1#4--#5#2#6}{ and~#3#1#4--#5#2#6}
  %   \crefname{hyp}{Hypothesis}{Hypotheses}
  %   
  %   \crefrangeformat{hyp}{Hypotheses~#3#1#4--#5#2#6}
  %   \crefmultiformat{hyp}{Hypotheses~#2#1#3}{ and~#2#1#3}{, #2#1#3}{ and~#2#1#3}
  %   \crefrangemultiformat{hyp}{Hypotheses~#3#1#4--#5#2#6}{ and~#3#1#4--#5#2#6}{, #3#1#4--#5#2#6}{ and~#3#1#4--#5#2#6}
  %   
  %   \crefname{rmk}{Remark}{Remarks}
  %   \crefrangeformat{rmk}{Remarks~#3#1#4--#5#2#6}
  %   \crefmultiformat{rmk}{Remarks~#2#1#3}{ and~#2#1#3}{, #2#1#3}{ and~#2#1#3}
  %   \crefrangemultiformat{rmk}{Remarks~#3#1#4--#5#2#6}{ and~#3#1#4--#5#2#6}{, #3#1#4--#5#2#6}{ and~#3#1#4--#5#2#6}
  %   
  %   % \newcommand\tikzfigure[5][0em]{%
  %   %   \begin{figure}[%
  %   %     pos=#2,%
  %   %     width=#4,%
  %   %     align=\centering%
  %   %   ]%
  %   %     \centering%
  %   %     \vspace*{#1}%
  %   %     \IfSubStr{#2}{p}{%
  %   %       \leavevmode\stepcounter{figposcount}%
  %   %       \zsaveposx{figpos-\thefigposcount}%
  %   %       \makebox[0pt][l]{%
  %   %         \hspace*{-\dimexpr\zposx{figpos-\thefigposcount}sp\relax}%
  %   %         \hspace*{\dimexpr0.5\paperwidth\relax}%
  %   %         \makebox[0pt][c]{%
  %   %           \scalebox{0.86}{\includegraphics{figs/#5.pdf}}%
  %   %         }%
  %   %       }
  %   %     }{%
  %   %       \includegraphics{figs/#5.pdf}%
  %   %     }%
  %   %     \caption{%
  %   %       \hyphenpenalty=0%
  %   %       \exhyphenpenalty=0%
  %   %       #3%
  %   %     }%
  %   %     \protect\label{#5}%
  %   %   \end{figure}%
  %   % }
  % %\endgroup
%\endgroup
